\chapter{Περιγραφή θέματος}

Το φαινόμενο του υπερτουρισμού, δηλαδή η διαταραχή της
ισορροπίας μεταξύ αριθμού επισκεπτών και αντοχής του
προορισμού από την τοπική κοινωνία και των υποδομών του,
απασχολεί την επιστημονική κοινότητα εδώ και δεκαετίες.
Τα κλασικά στατιστικά εργαλεία, όπως οι δείκτες τουριστικής
έντασης και πυκντότητας, συνήθως αδυνατούν να αναδείξουν την
δυναμική που δημιουργείται από την αλληλεπίδραση ετερογενών
πρακτόρων μέσα σε αστικό περιβάλλον. Για αυτό το λόγο, η
υπολογιστική μοντελοποίηση μέσω \en{Agent-Based Models (ABMs)},
όπως αναλύθηκε στην Ενότητα 2.2, κρίνεται κατάλληλη για την
μελέτη του υπερτουρισμού. Στο κεφάλαιο αυτό παρουσιάζεται μια
αντιπροσωπευτική εργασία που ακολουθεί αυτή τη προσέγγιση, καθώς
και ο στόχος και οι επεκτάσεις που εισάγει η παρούσα εργασία.

\section{Σχετικές εργασίες}

\subsection{\en{Moreno, Parada} και \en{Peña-Miranda}}

Οι \en{Moreno, Parada} και \en{Peña-Miranda} \cite{moreno2025agent}
αναπτύσσουν ένα από τα πληρέστερα \en{ABM} σε αστικό πλαίσιο
υπερτουρισμού, εφαρμοσμένο στην πόλη \en{Santa Marta} της Κολομβίας.
Η \en{Santa Marta} αποτελεί χαρακτηριστικό παράδειγμα μεσαίας πόλης
με ταχύτατα αναπτυσσόμενο τουρισμό, γεγονός που την καθιστά κατάλληλη
περίπτωση μελέτης για την ανάπτυξη ενός μοντέλου \en{ABM} που να
αναπαράγει τη δυναμική της τουριστικής πίεσης σε αστικό περιβάλλον.

Στο μοντέλο τους ορίζονται δύο κατηγορίες πρακτόρων: οι κάτοικοι, οι
οποίοι κινούνται για αναψυχή σε σημεία τουριστικού ενδιαφέροντος
\en{(SIT)} και επιστρέφουν στην κατοικία τους, και οι τουρίστες, οι
οποίοι εισέρχονται στο σύστημα από συγκεκριμένες εισόδους και κινούνται
σε οδικό δίκτυο βασισμένο σε γεωγραφικά δεδομένα \en{OpenStreetMap}.
Το λογισμικό που χρησιμοποιήθηκε για την υλοποίηση είναι το \en{GAMA}
πλαίσιο μοντελοποίησης πρακτόρων.
Οι συγγραφείς αξιολογούν τον κίνδυνο υπερτουρισμού μέσω κλασικών
δεικτών έντασης, πυκνότητας και συγκέντρωσης τουριστών. Προτείνουν
επίσης έναν νέο δείκτη τουριστικής έντασης, τον \en{Ism (Santa Maria
Overtourism Intensity)}, ο οποίος υπολογίζει μόνο το χρονικό διάστημα
που ο τουρίστας περνά πραγματικά στην πόλη, αποφεύγοντας την
υπερεκτίμηση της τουριστικής πίεσης που παρουσιάζουν οι κλασικοί
δείκτες.

Σε επίπεδο σεναρίων, δοκιμάζουν δύο στρατηγικές διαχείρισης
του υπερτουρισμού: την επέκταση των διαθέσιμων αξιοθέατων για χωρική
διασπορά των επισκεπτών και την παροχή πληροφοριών συνωστισμού στους
πράκτορες σε πραγματικό χρόνο. Και οι δύο στρατηγικές αποδεικνύονται
αποτελεσματικές στη μείωση της τουριστικής συγκέντρωσης,
καταδεικνύοντας την αξία της προσομοίωσης \en{ABM} ως εργαλείου
υποστήριξης αποφάσεων για διαχειριστές τουριστικών προορισμών.

Το μοντέλο αυτό, ωστόσο, εμφανίζει ορισμένες αδυναμίες. Οι τουρίστες
μοντελοποιούνται ως ομάδες με απλουστευμένη συμπεριφορά: επιλέγουν
τυχαία αξιοθέατα βάσει στατιστικών πιθανοτήτων, χωρίς να λαμβάνουν
υπόψη την ιεραρχία των σημείων ενδιαφέροντος, την οργανωτική
παρεκκλίνουσα συμπεριφορά εν κινήσει ή τα υποκειμενικά μεγέθη
εμπειρίας όπως ευχαρίστηση των επισκεπτών. Επιπλέον, το μοντέλο δε
μετρά ποιοτικές μεταβλητές αναφορικά με την εμπειρία των κατοίκων ή
των επισκεπτών ούτε παρέχει διαδραστικό εργαλείο για την
παρακολούθηση της προσομοίωσης σε πραγματικό χρόνο.

\subsection{\en{Yaprak, Okkiran} και \en{Vezali}}

Οι \en{Yaprak, Okkiran} και \en{Vezali} \cite{yaprak2025reframing}
προσεγγίζουν τον υπερτουρισμό από τη σκοπιά των μόνιμων κατοίκων,
εξετάζοντας πώς πέντε διαστάσεις βιωσιμότητας (οικονομική, περιβαλλοντική,
κοινωνικοπολιτισμική, πολιτική και τεχνολογική) διαμορφώνουν τις
αντιλήψεις τους για τον υπερτουρισμό και την υποστήριξή τους προς τον
τουρισμό. Η έρευνα εφαρμόζεται συγκριτικά σε δύο πόλεις της ανατολικής
Μεσογείου, την Αθήνα και την Κωνσταντινούπολη, που αντιμετωπίζουν ταχεία
ανάπτυξη τουριστικών ροών στα πλαίσια της μεταπανδημικής ανάκαμψης.
Η επιλογή αυτών των δύο πόλεων είναι ιδιαίτερα εύστοχη, διότι παρά τη
γεωγραφική τους εγγύτητα, διαφέρουν σημαντικά ως προς τη διοικητική δομή,
την αστική μορφολογία και τον τρόπο με τον οποίο απορροφούν τις
τουριστικές πιέσεις, καθιστώντας τες ιδανική βάση για συγκριτική ανάλυση.

Ως θεωρητικό πλαίσιο υιοθετείται το μοντέλο \en{Stimulus–Organism–Response
(S-O-R)}, όπου οι εξωτερικές πιέσεις  (οι πέντε διαστάσεις βιωσιμότητας)
λειτουργούν ως ερεθίσματα που διαμορφώνουν τις εσωτερικές αντιλήψεις των
κατοίκων, οι οποίες με τη σειρά τους καθορίζουν τη στάση τους απέναντι
στον τουρισμό. Τα δεδομένα συλλέχθηκαν μέσω ηλεκτρονικού ερωτηματολογίου
από 285 μόνιμους κατοίκους (136 από την Κωνσταντινούπολη και 149 από την
Αθήνα) και αναλύθηκαν με τη μέθοδο \en{PLS-SEM}. Για τη διασφάλιση της
συγκρισιμότητας μεταξύ των δύο δειγμάτων εφαρμόστηκε η διαδικασία
\en{MICOM} ελέγχου αναλλοίωτης μέτρησης.

Τα αποτελέσματα αναδεικνύουν σημαντικές διαφορές μεταξύ των δύο πόλεων.
Στην Αθήνα, η αυξημένη περιβαλλοντική, οικονομική και κοινωνικοπολιτισμική
συνειδητοποίηση ενισχύει αρνητικά τις αντιλήψεις υπερτουρισμού και μειώνει
τελικά την υποστήριξη των κατοίκων, ενώ στην Κωνσταντινούπολη η ίδια
αλληλουχία δε συνδέεται με στατιστικά σημαντική μείωση της υποστήριξης.
Οι διαστάσεις πολιτικής και τεχνολογικής βιωσιμότητας δεν παράγουν
στατιστικά σημαντικά αποτελέσματα σε καμία από τις δύο πόλεις, γεγονός που
οι συγγραφείς αποδίδουν στην περιορισμένη ορατότητα αυτών των εργαλείων
στο ευρύ κοινό.

Η εργασία συνεισφέρει σημαντικά στη βιβλιογραφία αναδεικνύοντας ότι οι
αντιλήψεις βιωσιμότητας λειτουργούν τοπικά και εξαρτώνται έντονα από το
αστικό, πολιτισμικό και διοικητικό πλαίσιο κάθε προορισμού, αμφισβητώντας
ενιαίες, καθολικές προσεγγίσεις διαχείρισης. Ωστόσο, η φύση της μεθοδολογίας,
δηλαδή εύκολη δειγματοληψία μέσω διαδικτύου και εγκάρσια σχεδίαση σε μια
χρονική στιγμή, δεν επιτρέπει αιτιώδεις συναρτήσεις, ενώ δε δίνει καμία
πληροφορία για τη χωρική ή χρονική διάσταση της τουριστικής πίεσης εντός της
πόλης.

\subsection{\en{Leka, Lagarias, Stratigea} και \en{Prekas}}

Οι \en{Leka, Lagarias, Stratigea} και \en{Prekas} \cite{leka2025methodological}
αναπτύσσουν ένα ολοκληρωμένο μεθοδολογικό πλαίσιο αξιολόγησης υπερτουρισμού
με έμφαση στη γεωχωρική διάσταση του φαινομένου, εφαρμοσμένο στη Σαντορίνη,
έναν από τους πλέον φορτωμένους τουριστικά προορισμούς της Μεσογείου. Η
επιλογή της Σαντορίνης ως πεδίου εφαρμογής δεν είναι τυχαία. Το νησί συνδυάζει
γεωγραφική απομόνωση, μονοκαλλιεργητικό μοντέλο τουρισμού και έντονες
πιέσεις σε φυσικούς και πολιτιστικούς πόρους, καθιστώντας το χαρακτηριστικό
παράδειγμα νησιωτικού υπερτουρισμού.

Το πλαίσιο βασίζεται στο μοντέλο \en{PSR (Pressure/State/Response)} και
συνδυάζει δεδομένα τουριστικής προσφοράς και ζήτησης με περιβαλλοντικούς,
οικονομικούς και κοινωνικούς δείκτες. Οι δείκτες οργανώνονται σε δύο βασικές
κατηγορίες: δείκτες τουριστικής έντασης, που εκφράζουν την τουριστική πίεση
ανά κάτοικο, και δείκτες τουριστικής πυκνότητας, που εκφράζουν την πίεση ανά
μονάδα έκτασης. Η γεωχωρική ανάλυση πραγματοποιείται μέσω δορυφορικών
εικόνων \en{Sentinel-2} και βάσεων δεδομένων \en{Copernicus}, επιτρέποντας
τον εντοπισμό και την ποσοτικοποίηση χωρικών επιπτώσεων όπως η αστική
εξάπλωση, η υποβάθμιση βλάστησης και η απώλεια αγροτικής γης. Το
μεθοδολογικό πλαίσιο σχεδιάζεται ρητά ώστε να είναι επαναλήψιμο σε
οποιονδήποτε τουριστικό προορισμό, αξιοποιώντας παγκοσμίως διαθέσιμες βάσεις
δεδομένων ανοιχτής πρόσβασης.

Τα αποτελέσματα αποκαλύπτουν ότι η Σαντορίνη έχει ξεπεράσει κρίσιμα όρια
φέρουσας ικανότητας σε πολλαπλές διαστάσεις. Η τουριστική υποδομή εισβάλλει
σε αρχαιολογικούς χώρους και ευάλωτα τμήματα του νησιού, ενώ η εξάπλωση
\en{Airbnb} και ξενοδοχειακών μονάδων προκαλεί αλλαγές χρήσης γης σε ζώνες
\en{Natura} 2000. Παράλληλα, το κόστος ενοικίασης κατοικίας υπερβαίνει κατά
56\% τον μέσο όρο της ευρύτερης περιοχής, ενώ οι τιμές πώλησης είναι
διπλάσιες, γεγονός που δυσχεραίνει σοβαρά τη ζωή των μόνιμων κατοίκων.
Οι συγγραφείς εντοπίζουν επίσης χωρικές ανισότητες εντός του νησιού, με
ορισμένες ζώνες, όπως το κεντρικό τμήμα και η Οία, να φέρουν πολύ μεγαλύτερο
φορτίο από άλλες.

Παρά τη σημαντική συνεισφορά του, το πλαίσιο αυτό αξιολογεί εκ των υστέρων
ήδη καταγεγραμμένα δεδομένα, χωρίς να προσφέρει τη δυνατότητα δυναμικής
προσομοίωσης σεναρίων. Επιπλέον, η προσέγγιση είναι αποκλειστικά
ποσοτική-χαρτογραφική, καθώς δεν μοντελοποιεί τη συμπεριφορά μεμονωμένων
πρακτόρων. Τέλος, η εφαρμογή σε νησιωτικό προορισμό θέτει ορισμένες
μεθοδολογικές ιδιαιτερότητες που δεν μεταφέρονται αυτόματα σε αστικά
περιβάλλοντα ηπειρωτικών πόλεων.


\section{Στόχος και Προσέγγιση της Παρούσας Εργασίας}

Η παρούσα εργασία επεκτείνει τις παραπάνω προσέγγισεις σε τρεις άξονες.
Πρώτων, εισάγει πλουσιότερο γνωστικό μοντέλο συμπεριφοράς πρακτόρων:
οι τουρίστες δεν επιλέγουν τυχαία προορισμό, αλλά ακολουθούν μια
ιεράρχημένη λογική προτεραιοτήτων τριών επιπέδων (\en{Tier A}:
αξιοθέατα υψηλής αξίας, \en{Tier B}: καφέ και εστιατόρια, \en{Tier C}:
καταστήματα). Προβλέπεται επίσης μηχανισμός εμπρόσθιας οπτικής σάρωσης
\en{(forward vision scan)} σε κώνο 30 μοιρών, ο οποίος επιτρέπει στον
τουρίστα να εντοπίζει και να επισκεφτεί ευκαιριακά σταθμούς του
\en{Tier B} ή \en{C} καθώς μετακινείται προς τον κυρίως στόχο του.

Δεύτερον, ενσωματώνονται υποκειμενικές μεταβλητές που αντικατοπτρίζουν
την ανθρώπινη διάσταση του φαινομένου: η ευτυχία των κατοίκων
\en{(happiness)} ως συνάρτηση της τουριστικής πίεσης στην γειτονιά
τους, και η ικανοποίηση των τουριστών \en{(satisfaction)} ως συνάρτηση
του επιπέδου συνωστισμού στους χώρους που επισκέπτονται. Αυτές οι
μεταβλητές αναδεικνύουν την ποσότητα της αλληλεπίδρασης μεταξύ
κατοίκων και τουριστών, κάτι που δεν καλύπτεται από κανένα από τους
κλασικούς δείκτες της βιβλιογραφίας.

Τρίτον, η περιοχή μελέτης είναι η Πάτρα, μια αστική πόλη μεσαίου
μεγέθους και ιδιαίτερα ανεπτυγμένη τουριστική υποδομή στην ιστορική
της περιοχή. Η πόλη διαθέτει αξιοθέατα πολλαπλών κατηγοριών - ιστορικά
μνημεία, μουσεία, χώροι εστίασης και καταστήματα - καθιστώντας ιδανική
περίπτωση για την εφαρμογή ιεραρχίας \en{POI} τριών επιπέδων. Τα
γεωγραφικά δεδομένα προέρχονται από \en{OpenStreetMap} μέσω του
αρχείου δεδομένων \en{Greece Shapefile}, ενσωματώνοντας το οδικό δίκτυο
και τα σημεία ενδιαφέροντος της πόλης.

Τέλος, η υλοποίηση γίνεται σε περιβάλλον ανοιχτού κώδικα με \en{Python}
και το πλαίσιο \en{Mesa}, ενώ η οπτικοποίηση γίνεται μέσω διαδραστικού
χάρτη βασισμένο στη βιβλιοθήκη \en{Leaflet} της \en{Javascript},
παρέχοντας στον χρήστη τη δυνατότητα παρακολούθησης της προσομοίωσης σε
πραγματικό χρόνο, χαρτογράφησης θερμικών χαρτών θορύβου και ρύπανσης
και γραφημάτων κατανομής ευτυχίας και ικανοποίησης. Αυτή η προσέγγιση
επιτρέπει τη διαδραστική εξερεύνηση των αναδυόμενων σχεδίων τουριστικής
κίνησης και κατανομής σε σχέση με την ευτθχία των κατοίκων.
